%% GENERATED by tools/from_docs.py from stepss-docs. Do not edit by hand:
%% edit the page named below in stepss-docs and regenerate.
%% source: stepss-docs/src/content/docs/models/custom-exciters.mdx

\chapter{Other excitation controllers}\label{doc:model-exc-custom}

This page documents the custom excitation system models available in RAMSES. These models range from a trivial constant-voltage source to full-featured AVR with PSS and OEL. Models without a \texttt{.txt} definition file are implemented directly in Fortran and are described from their source headers.

\begin{docsnote}{Usage in Dynamic Data Files}

Exciter models in RAMSES are defined as sub-records within a \texttt{SYNC\_MACH} block using the \texttt{EXC} keyword, followed by the model name and its parameters. They are not standalone records. The \texttt{EXC} line appears after the machine parameters and before the \texttt{TOR} (torque/governor) line. The entire block ends with a single \texttt{;} after the \texttt{TOR} record.

\begin{Verbatim}
SYNC_MACH name bus FP FQ P Q SNOM Pnom H D IBRATIO
                    XT  Xl  Xd  X'd  X"d  Xq  X'q  X"q  m  n  Ra  T'do  T"do  T'qo  T"qo
              EXC model_name  param1  param2  ...
              TOR model_name  param1  param2  ... ;
\end{Verbatim}

This page documents the exciters that are not IEEE-standard models: \texttt{CONSTANT}, \texttt{1ST\_ORDER}, \texttt{GENERIC1} and \texttt{GENERIC2} (built in, uppercase, no prefix), plus \texttt{kundur}, \texttt{GENERIC} and \texttt{AVR\_DG} (registered, and RAMSES adds the \texttt{exc\_} prefix automatically).

The IEEE-standard exciters are on the \hyperref[doc:model-exc-ieee]{IEEE Exciter Models} page. For the whole catalogue at a glance, including which models are compiled but not callable by name, see the \hyperref[doc:model-index]{Model Reference index}.

\end{docsnote}

\begin{center}\rule{0.5\linewidth}{0.5pt}\end{center}

\section{\texorpdfstring{CONSTANT (\texttt{exc\_constant})}{CONSTANT (exc\_constant)}}\label{model-exc-custom:constant-exc_constant}

The data-file model name is \texttt{CONSTANT}.

\subsection{Scientific Description}\label{model-exc-custom:scientific-description}

The \texttt{exc\_constant} model provides a \textbf{constant field voltage} source. The field voltage \(V_f\) is fixed at its load-flow initialisation value and does not respond to any input:

\[
V_f(t) = V_{f0} \quad \forall t
\]

The single algebraic state satisfies:

\[
f(1) = x(1) - V_{f0} = 0
\]

This model is useful for open-loop tests, for machines that are not equipped with a voltage regulator, or as a placeholder during model development.

\textbf{Implementation details (from Fortran source):}
- \texttt{nbxvar\ =\ 1}, \texttt{nbzvar\ =\ 0}, \texttt{nbdata\ =\ 0}
- One additional parameter: \texttt{prm(1)\ =\ Vf0} (set at initialisation from the load-flow field voltage)
- Observables: \texttt{vf}, \texttt{if}

\subsection{Parameters}\label{model-exc-custom:parameters}

\begin{small}\begin{longtable}[]{@{}
  >{\raggedright\arraybackslash}p{(\linewidth - 4\tabcolsep) * \real{0.2550}}
  >{\raggedright\arraybackslash}p{(\linewidth - 4\tabcolsep) * \real{0.7450}}
@{}}
\toprule\noalign{}
\begin{minipage}[b]{\linewidth}\raggedright
Parameter
\end{minipage} & \begin{minipage}[b]{\linewidth}\raggedright
Description
\end{minipage} \\
\midrule\noalign{}
\endhead
\bottomrule\noalign{}
\endlastfoot
\emph{(none)} & No user-supplied data parameters. The field voltage is taken from the initialisation (load-flow) value automatically. \\
\end{longtable}\end{small}

\subsection{Usage Example}\label{model-exc-custom:usage-example}

\textbf{RAMSES Data File}

\begin{Verbatim}
SYNC_MACH GEN1  BUS1  1. 1.  0.  0.   600. 570. 3. 0. 0.95
                    XT   0.15  1.1  0.25  0.2  0.7  *  0.2  0.1  6.0257  0.  5.00  0.05  *  0.1
              EXC CONSTANT
              TOR CONSTANT ;
\end{Verbatim}

\textbf{Notes}

\begin{itemize}
\tightlist
\item
  No parameters need to be supplied; the model self-initialises to the load-flow \(V_f\) value.
\item
  Useful for simulation of machines operating in manual excitation mode.
\item
  The model name in the data file is \texttt{CONSTANT} (not \texttt{exc\_constant}).
\end{itemize}

\begin{center}\rule{0.5\linewidth}{0.5pt}\end{center}

\section{\texorpdfstring{1ST\_ORDER (\texttt{exc\_1storder})}{1ST\_ORDER (exc\_1storder)}}\label{model-exc-custom:st_order-exc_1storder}

The data-file model name is \texttt{1ST\_ORDER} (uppercase, with the underscore).

\subsection{Scientific Description}\label{model-exc-custom:scientific-description-1}

The \texttt{exc\_1storder} model is a \textbf{first-order automatic voltage regulator} (AVR). It represents the simplest dynamic exciter: a proportional gain \(K_A\) acting on the terminal voltage error, filtered through a single first-order lag \(T_A\), with output limits.

The AVR computes the voltage error and drives the field voltage:

\[
\frac{d V_f}{dt} = \frac{1}{T_A}\left[K_A\left(V_{ref} - V\right) - V_f\right], \quad V_f \in [V_{f,min},\, V_{f,max}]
\]

The reference voltage \(V_{ref}\) is initialised from:

\[
V_{ref} = V_0 + \frac{V_{f0}}{K_A}
\]

\textbf{Implementation details (from Fortran source):}
- \texttt{nbxvar\ =\ 1}, \texttt{nbzvar\ =\ 1}, \texttt{nbdata\ =\ 4}
- Parameters: \texttt{prm(1)\ =\ KA}, \texttt{prm(2)\ =\ TA}, \texttt{prm(3)\ =\ Vfmin}, \texttt{prm(4)\ =\ Vfmax}
- Additional parameter: \texttt{prm(5)\ =\ V0} (voltage setpoint, computed at initialisation)
- The ODE in the continuous mode is:

\[
f(1) = \frac{-V_f + K_A(V_{ref} - V)}{T_A}
\]

\begin{itemize}
\tightlist
\item
  With anti-windup: if the limiter is active, \texttt{f(1)\ =\ Vf\ -\ Vfmin} (or \texttt{Vfmax}).
\end{itemize}

\subsection{Parameters}\label{model-exc-custom:parameters-1}

\begin{small}\begin{longtable}[]{@{}ll@{}}
\toprule\noalign{}
Parameter & Description \\
\midrule\noalign{}
\endhead
\bottomrule\noalign{}
\endlastfoot
\texttt{KA} & AVR gain (pu/pu) \\
\texttt{TA} & AVR time constant (s) \\
\texttt{Vfmin} & Minimum field voltage (pu) \\
\texttt{Vfmax} & Maximum field voltage (pu) \\
\end{longtable}\end{small}

\subsection{Usage Example}\label{model-exc-custom:usage-example-1}

\textbf{RAMSES Data File}

\begin{Verbatim}
SYNC_MACH GEN1  BUS1  1. 1.  0.  0.   600. 570. 3. 0. 0.95
                    XT   0.15  1.1  0.25  0.2  0.7  *  0.2  0.1  6.0257  0.  5.00  0.05  *  0.1
              EXC 1ST_ORDER  200.0  0.05  0.0  6.0
                            ! KA    TA    Vfmin Vfmax
              TOR CONSTANT ;
\end{Verbatim}

\textbf{Notes}

\begin{itemize}
\tightlist
\item
  This model is the Fortran-coded equivalent of a simple \texttt{tf1plim} block with input \(V_{ref} - V\).
\item
  \texttt{KA} and \texttt{TA} together set the closed-loop bandwidth.
\item
  The model name in the data file is \texttt{1ST\_ORDER} (uppercase, with the underscore).
\end{itemize}

\begin{center}\rule{0.5\linewidth}{0.5pt}\end{center}

\section{\texorpdfstring{kundur (\texttt{exc\_kundur})}{kundur (exc\_kundur)}}\label{model-exc-custom:kundur-exc_kundur}

The data-file model name is \texttt{kundur} or \texttt{exc\_kundur}. The \texttt{exc\_} prefix is added automatically.

\subsection{Scientific Description}\label{model-exc-custom:scientific-description-2}

The \texttt{exc\_kundur} model implements the \textbf{excitation system from Kundur's textbook} (\emph{Power System Stability and Control}, 1994). It includes a first-order voltage transducer, a proportional AVR with transient gain reduction (TGR), and a two-stage PSS acting on rotor speed deviation.

\textbf{Voltage transducer:}

\[
\frac{d V_{fil}}{dt} = \frac{1}{T_R}(V - V_{fil})
\]

\textbf{AVR with transient gain reduction (lead--lag):}

\[
V_f(s) = K_A \cdot \frac{1 + s T_A}{1 + s T_B} \cdot \left(V_{ref} - V_{fil} + V_s\right)
\]

where \(V_{ref} = V + V_f / K_A\) is computed at initialisation and \(V_s\) is the PSS output.

\textbf{PSS (two-stage lead--lag with washout):}

\[
x_{wash}(s) = \frac{K_{STAB} \cdot s T_W}{1 + s T_W} \cdot \omega(s)
\]

\[
V_s(s) = \frac{1 + s T_1}{1 + s T_2} \cdot \frac{1 + s T_3}{1 + s T_4} \cdot x_{wash}(s)
\]

This is the standard two-stage phase-compensating PSS structure. The CODEGEN blocks in the \texttt{.txt} file are:

\begin{small}\begin{longtable}[]{@{}lll@{}}
\toprule\noalign{}
Block & CODEGEN type & Function \\
\midrule\noalign{}
\endhead
\bottomrule\noalign{}
\endlastfoot
Voltage filter & \texttt{tf1p} & \(1/(1+sT_R)\) \\
Main summing junction & \texttt{algeq} & \(V_{ref} - V_{fil} + V_s - dv\) \\
AVR (TGR) & \texttt{tf1p1z} & \(K_A(1+sT_A)/(1+sT_B)\) \\
PSS washout & \texttt{tfder1p} & \(K_{STAB} \cdot sT_W/(1+sT_W)\) \\
PSS lead--lag 1 & \texttt{tf1p1z} & \((1+sT_1)/(1+sT_2)\) \\
PSS lead--lag 2 & \texttt{tf1p1z} & \((1+sT_3)/(1+sT_4)\) \\
\end{longtable}\end{small}

\subsection{Parameters}\label{model-exc-custom:parameters-2}

\begin{small}\begin{longtable}[]{@{}ll@{}}
\toprule\noalign{}
Parameter & Description \\
\midrule\noalign{}
\endhead
\bottomrule\noalign{}
\endlastfoot
\texttt{TR} & Voltage measurement filter time constant (s) \\
\texttt{KA} & AVR gain (pu/pu) \\
\texttt{TA} & TGR lead time constant (s) \\
\texttt{TB} & TGR lag time constant (s) \\
\texttt{KSTAB} & PSS gain \\
\texttt{TW} & PSS washout time constant (s) \\
\texttt{T1} & PSS first lead--lag zero time constant (s) \\
\texttt{T2} & PSS first lead--lag pole time constant (s) \\
\texttt{T3} & PSS second lead--lag zero time constant (s) \\
\texttt{T4} & PSS second lead--lag pole time constant (s) \\
\end{longtable}\end{small}

\textbf{Computed at initialisation:} \texttt{Vref\ =\ V\ +\ Vf\ /\ KA}

\subsection{Usage Example}\label{model-exc-custom:usage-example-2}

\textbf{RAMSES Data File}

\begin{Verbatim}
SYNC_MACH GEN1  BUS1  1. 1.  0.  0.   600. 570. 3. 0. 0.95
                    XT   0.15  1.1  0.25  0.2  0.7  *  0.2  0.1  6.0257  0.  5.00  0.05  *  0.1
              EXC kundur  0.02  200.0  0.02  10.0  9.5  1.41  0.154  0.033  0.154  0.033
                          ! TR   KA     TA    TB    KSTAB TW    T1     T2     T3     T4
              TOR CONSTANT ;
\end{Verbatim}

\textbf{Notes}

\begin{itemize}
\tightlist
\item
  Parameters match the Example 12.6 from Kundur (1994), applicable to a medium-speed generator.
\item
  The PSS stabilises local modes in the 0.5--2 Hz range.
\item
  No output limiter is defined in the base model; add clipping externally if required.
\item
  The model name in the data file is \texttt{kundur} (not \texttt{exc\_kundur}).
\end{itemize}

\begin{center}\rule{0.5\linewidth}{0.5pt}\end{center}

\section{\texorpdfstring{GENERIC1 (\texttt{exc\_generic1})}{GENERIC1 (exc\_generic1)}}\label{model-exc-custom:generic1-exc_generic1}

The data-file model name is \texttt{GENERIC1}.

\subsection{Scientific Description}\label{model-exc-custom:scientific-description-3}

The \texttt{exc\_generic1} model is a \textbf{generic AVR with OEL (Over-Excitation Limiter) and PSS}, implemented in Fortran. It includes a proportional--integral AVR, a PSS acting on rotor speed or active power, and an inverse-time OEL.

\textbf{AVR:}

The excitation system uses a first-order integrating controller (washout or PI-type, set by the gain \(G\) and time constants \(T_A\), \(T_B\)) with output limits \([L_3,\, L_4]\):

\[
V_f(s) = G \cdot \frac{1 + s T_A}{1 + s T_B} \cdot \left(V_{ref} - V\right) \cdot \frac{1}{1 + s T_E}
\]

\textbf{PSS (speed or power input, selectable):}

Input signal is selected by \texttt{SPEEDIN}: 1 = rotor speed \(\omega\), 0 = electrical power \(P\).

\[
V_{PSS}(s) = K_{PSS} \cdot \frac{s T_W}{1 + s T_W} \cdot \frac{1 + s T_1}{1 + s T_2} \cdot \frac{1 + s T_3}{1 + s T_4} \cdot u(s)
\]

Output is clamped to \([DV_{MIN},\, DV_{MAX}]\).

\textbf{OEL (inverse-time limiter):}

The OEL monitors field current \(I_f\). When \(I_f > IFLIM\), a timer begins accumulating via:

\begin{itemize}
\tightlist
\item
  Dead zone \(d\) and phase \(f\) parameters control OEL activation thresholds.
\item
  The OEL integrates the overcurrent error and reduces \(V_{ref}\) once the timer reaches threshold \(L_1\), with output range \([L_1, L_2]\).
\end{itemize}

\textbf{Parameters from Fortran source} (\texttt{prm} array):

\begin{small}\begin{longtable}[]{@{}lll@{}}
\toprule\noalign{}
Index & Name & Description \\
\midrule\noalign{}
\endhead
\bottomrule\noalign{}
\endlastfoot
1 & \texttt{IFLIM} & Field current OEL threshold (pu) \\
2 & \texttt{d} & OEL dead zone (pu) \\
3 & \texttt{f} & OEL phase characteristic \\
4 & \texttt{S} & AVR setpoint input (usually 0 = \(V\)) \\
5 & \texttt{K1} & AVR lead numerator gain \\
6 & \texttt{K2} & AVR lag denominator gain \\
7 & \texttt{L1} & OEL integrator lower limit (pu) \\
8 & \texttt{L2} & OEL integrator upper limit (pu) \\
9 & \texttt{G} & AVR overall gain \\
10 & \texttt{TA} & AVR lead time constant (s) \\
11 & \texttt{TB} & AVR lag time constant (s) \\
12 & \texttt{TE} & Exciter time constant (s) \\
13 & \texttt{L3} & AVR output lower limit (pu) \\
14 & \texttt{L4} & AVR output upper limit (pu) \\
15 & \texttt{SPEEDIN} & PSS input signal: 1 = speed, 0 = power \\
16 & \texttt{KPSS} & PSS gain \\
17 & \texttt{Tw} & PSS washout time constant (s) \\
18 & \texttt{T1} & PSS lead--lag 1 zero (s) \\
19 & \texttt{T2} & PSS lead--lag 1 pole (s) \\
20 & \texttt{T3} & PSS lead--lag 2 zero (s) \\
21 & \texttt{T4} & PSS lead--lag 2 pole (s) \\
22 & \texttt{DVMIN} & PSS output lower limit (pu) \\
23 & \texttt{DVMAX} & PSS output upper limit (pu) \\
\end{longtable}\end{small}

\subsection{Usage Example}\label{model-exc-custom:usage-example-3}

\textbf{RAMSES Data File}

\begin{Verbatim}
SYNC_MACH GEN1  BUS1  1. 1.  0.  0.   600. 570. 3. 0. 0.95
                    XT   0.15  1.1  0.25  0.2  0.7  *  0.2  0.1  6.0257  0.  5.00  0.05  *  0.1
              EXC GENERIC1  1.05  -0.1  0.5  0.0  1.0  1.0  0.0  0.5  200.0  0.02  0.10  0.05  0.0  6.0
                            ! IFLIM  d    f    S   K1   K2   L1   L2   G      TA    TB    TE    L3   L4
                            1     9.5   1.40  0.15  0.03  0.15  0.03  -0.10  0.10
                            ! SPEEDIN KPSS  Tw    T1    T2    T3    T4    DVMIN  DVMAX
              TOR CONSTANT ;
\end{Verbatim}

\begin{center}\rule{0.5\linewidth}{0.5pt}\end{center}

\section{\texorpdfstring{AVR\_DG (\texttt{exc\_AVR\_DG})}{AVR\_DG (exc\_AVR\_DG)}}\label{model-exc-custom:avr_dg-exc_avr_dg}

The data-file model name is \texttt{AVR\_DG}. Registered since 3.76.

\subsection{Scientific Description}\label{model-exc-custom:scientific-description-4}

An AVR for \textbf{dispersed generation}: the \hyperref[model-exc-custom:generic1-exc_generic1]{\texttt{GENERIC1}} excitation
system, OEL and PSS, plus a \textbf{reactive-power PI controller} that trims the voltage
reference so the unit holds a reactive-power setpoint.

Parameters 1 to 23 are those of \texttt{GENERIC1}, in the same positions and with the same
meanings, and the AVR, PSS and OEL behave identically. Only the summing junction changes,
gaining a third contribution:

\[
e = V_{ref} - V + \Delta V_{PSS} + \Delta V_{Q}
\]

\textbf{Reactive-power regulator:}

The error is the deviation from the setpoint captured at initialisation:

\[
\dot{x}_{Q} = K_{QI} \left(Q^{*} - Q\right), \qquad
\Delta V_{Q} = K_{QP} \left(Q^{*} - Q\right) + x_{Q}
\]

clamped to \([L_5,\, L_6]\). Setting \(K_{QP} = K_{QI} = 0\) leaves the loop dormant and the
model reduces to \texttt{GENERIC1}.

\textbf{Parameters:}

Indices 1 to 23 are exactly those of \hyperref[model-exc-custom:generic1-exc_generic1]{\texttt{GENERIC1}}. The four that
follow are new:

\begin{small}\begin{longtable}[]{@{}lll@{}}
\toprule\noalign{}
Index & Name & Description \\
\midrule\noalign{}
\endhead
\bottomrule\noalign{}
\endlastfoot
24 & \texttt{KQP} & Reactive-power regulator proportional gain \\
25 & \texttt{KQI} & Reactive-power regulator integral gain \\
26 & \texttt{L5} & Regulator output lower limit (pu) \\
27 & \texttt{L6} & Regulator output upper limit (pu) \\
\end{longtable}\end{small}

\textbf{Implementation details:}
- \texttt{nbxvar\ =\ 8}, \texttt{nbzvar\ =\ 7}, \texttt{nbdata\ =\ 27}, \texttt{nbaddpar\ =\ 2}
- Additional: \texttt{prm(28)\ =\ V0} (voltage reference) and \texttt{prm(29)\ =\ Qset}, both captured at
initialisation
- Observables: those of \texttt{GENERIC1} plus \texttt{deltavq} and \texttt{Qset}

\subsection{Usage Example}\label{model-exc-custom:usage-example-4}

\textbf{RAMSES Data File}

\begin{Verbatim}
SYNC_MACH G4  N4 1. 1.  0.  0. 5 4.5 3. 0. 1.9333
              XT  0.0937  2.027  0.210  0.173  2.027  *  0.173  0.1  6  0.00  3.1  0.0387  *  0.0387
              EXC AVR_DG  2.865 -0.1  0.  1.  100. -1. -20  10  50  4  20  0.1  0  5
                          ! IFLIM  d    fc   S   K1   K2   L1  L2  G  TA  TB  TE  L3 L4
                          1  0.  5.  1.  1.  1.  1.  0.  0.
                          ! SPEEDIN KPSS Tw  T1  T2  T3  T4  DVMIN DVMAX
                          0.04  0.06  -0.2  0.33
                          ! KQP   KQI    L5    L6
              TOR HYDRO_DG  0.0  0.1  0.4  0.2  0.1  1.0 ;
\end{Verbatim}

\textbf{Notes}

\begin{itemize}
\tightlist
\item
  The setpoint is \textbf{not} a parameter. \texttt{Qset} is captured from the power-flow solution at
  initialisation, so the unit holds whatever reactive power it was dispatched at.
\item
  With \texttt{KQP} and \texttt{KQI} non-zero the unit regulates reactive power rather than terminal
  voltage, which is the usual arrangement for dispersed generation connected to a
  distribution network under a connection agreement.
\item
  \texttt{KPSS\ =\ 0} disables the stabiliser. Note that this drives the PSS branch to its lower
  limit, so \texttt{DVMIN} should be 0 as well, as in the example above.
\item
  Pair it with \hyperref[doc:model-tor-custom]{\texttt{HYDRO\_DG}} for a unit that
  regulates both its active and reactive power and provides no primary frequency response.
\end{itemize}

\begin{center}\rule{0.5\linewidth}{0.5pt}\end{center}

\section{\texorpdfstring{GENERIC2 (\texttt{exc\_generic2})}{GENERIC2 (exc\_generic2)}}\label{model-exc-custom:generic2-exc_generic2}

The data-file model name is \texttt{GENERIC2}.

\subsection{Scientific Description}\label{model-exc-custom:scientific-description-5}

The \texttt{exc\_generic2} model is a more detailed \textbf{generic AVR}. It includes a voltage compensator, a combined PI--proportional regulator, an AC exciter with saturation, a PSS, and an OEL.

\textbf{Voltage measurement with line-drop compensation:}

\[
V_c = V + X_c \cdot I_q - R_c \cdot I_d \quad (\text{or simplified})
\]

\textbf{PI--proportional AVR with transient gain reduction:}

\[
V_{AVR}(s) = \left(K_P + \frac{K_I}{s} + C \cdot \frac{1 + s T_C}{1 + s T_B}\right) \cdot K_A \cdot \frac{1}{1 + s T_A}(V_{ref} - V_c)
\]

\textbf{AC exciter with saturation:}

\[
\frac{d V_E}{dt} = \frac{1}{T_E}\left[V_R - K_E \cdot V_E - S_E(V_E)\right]
\]

where \(S_E(V_E)\) is the piecewise-linear saturation function defined by \((E_1, S_{E1})\) and \((E_2, S_{E2})\).

\textbf{PSS (type selectable by \texttt{typss}):}

\[
V_s(s) = K_{PSS} \cdot \frac{s T_W}{1 + s T_W} \cdot \frac{1 + s T_1}{1 + s T_2} \cdot \frac{1 + s T_3}{1 + s T_4} \cdot u(s)
\]

\textbf{OEL (inverse-time, three-segment):}

The OEL uses a three-level piecewise-linear characteristic \((s_1, T_{oel1})\), \((s_2, T_{oel2})\), \((s_3, T_{oel3})\) to compute the allowable overcurrent time.

\textbf{Parameters from Fortran \texttt{associate} block} (39 data parameters + 4 additional):

\begin{small}\begin{longtable}[]{@{}lll@{}}
\toprule\noalign{}
Index & Name & Description \\
\midrule\noalign{}
\endhead
\bottomrule\noalign{}
\endlastfoot
1 & \texttt{Xc} & Line-drop compensation reactance (pu) \\
2 & \texttt{Tr} & Voltage measurement time constant (s) \\
3 & \texttt{C} & TGR numerator gain \\
4 & \texttt{Tc} & TGR lead time constant (s) \\
5 & \texttt{Tb} & TGR lag time constant (s) \\
6 & \texttt{Ka} & AVR proportional gain \\
7 & \texttt{Ta} & AVR time constant (s) \\
8 & \texttt{Ki} & AVR integral gain \\
9 & \texttt{Kp} & AVR proportional path gain \\
10 & \texttt{VrmaxD} & Regulator max (dynamic, pu) \\
11 & \texttt{VrminD} & Regulator min (dynamic, pu) \\
12 & \texttt{G} & Exciter overall gain \\
13 & \texttt{Ke} & Exciter self-excitation coefficient \\
14 & \texttt{Te} & Exciter time constant (s) \\
15 & \texttt{E1} & Saturation breakpoint 1 (pu) \\
16 & \texttt{SeE1} & Saturation factor at \(E_1\) \\
17 & \texttt{E2} & Saturation breakpoint 2 (pu) \\
18 & \texttt{SeE2} & Saturation factor at \(E_2\) \\
19 & \texttt{Vrmax} & Field voltage upper limit (pu) \\
20 & \texttt{Vrmin} & Field voltage lower limit (pu) \\
21 & \texttt{typss} & PSS input type (1 = speed, 2 = power) \\
22 & \texttt{W} & PSS input gain \\
23 & \texttt{Tf} & PSS derivative filter time constant (s) \\
24 & \texttt{Kpss} & PSS gain \\
25 & \texttt{Tw} & PSS washout time constant (s) \\
26 & \texttt{T1} & PSS lead--lag 1 zero (s) \\
27 & \texttt{T2} & PSS lead--lag 1 pole (s) \\
28 & \texttt{T3} & PSS lead--lag 2 zero (s) \\
29 & \texttt{T4} & PSS lead--lag 2 pole (s) \\
30 & \texttt{Lim} & PSS output symmetric limit (pu) \\
31 & \texttt{LimSup} & PSS upper asymmetric limit (pu) \\
32 & \texttt{s1} & OEL segment 1 (pu \(I_f\)) \\
33 & \texttt{toel1} & OEL time at segment 1 (s) \\
34 & \texttt{s2} & OEL segment 2 (pu \(I_f\)) \\
35 & \texttt{toel2} & OEL time at segment 2 (s) \\
36 & \texttt{s3} & OEL segment 3 (pu \(I_f\)) \\
37 & \texttt{toel3} & OEL time at segment 3 (s) \\
38 & \texttt{perm} & OEL permanent current limit (pu) \\
39 & \texttt{LimInf} & PSS lower asymmetric limit (pu) \\
\end{longtable}\end{small}

\subsection{Usage Example}\label{model-exc-custom:usage-example-5}

\textbf{RAMSES Data File}

\begin{Verbatim}
SYNC_MACH GEN1  BUS1  1. 1.  0.  0.   600. 570. 3. 0. 0.95
                    XT   0.15  1.1  0.25  0.2  0.7  *  0.2  0.1  6.0257  0.  5.00  0.05  *  0.1
              EXC GENERIC2  0.0  0.02  1.0  0.10  1.00  200.0  0.05  5.0  1.0  6.0  -6.0  1.0  0.0  0.50
                            ! Xc   Tr   C    Tc    Tb    Ka     Ta    Ki   Kp  VrmaxD VrminD G   Ke  Te
                            3.10  0.066  4.18  0.91  6.0  -6.0  1.0  1.0  0.05  9.5  1.40  0.154  0.033  0.154  0.033
                            ! E1   SeE1   E2   SeE2  Vrmax Vrmin typss W   Tf   Kpss  Tw    T1     T2     T3     T4
                            0.10  0.10  1.1  200.  1.2  60.  1.5  10.  1.05  -0.10
                            ! Lim  LimSup s1  toel1 s2  toel2 s3  toel3 perm  LimInf
              TOR CONSTANT ;
\end{Verbatim}

\begin{center}\rule{0.5\linewidth}{0.5pt}\end{center}

\section{\texorpdfstring{GENERIC (\texttt{exc\_GENERIC})}{GENERIC (exc\_GENERIC)}}\label{model-exc-custom:generic-exc_generic}

The data-file model name is \texttt{exc\_GENERIC} (or \texttt{GENERIC}. The \texttt{exc\_} prefix is added automatically).

\begin{docsnote}{Renamed in RAMSES 3.57}

This model was called \texttt{exc\_GENERIC3} up to and including 3.56, and a second
registered name, \texttt{exc\_GENERIC4}, resolved to a byte-for-byte copy of it. Both
names were removed in 3.57 in favour of the single name \texttt{exc\_GENERIC}.

On 3.57 and later, a record naming \texttt{exc\_GENERIC3} or \texttt{exc\_GENERIC4} stops with
\texttt{Excitation\ model\ not\ found\ in\ definition}. Change it to \texttt{exc\_GENERIC}; the 20
data parameters and their order are unchanged, so nothing else in the record
moves.

\end{docsnote}

\subsection{Scientific Description}\label{model-exc-custom:scientific-description-6}

The \texttt{exc\_GENERIC} model (note capital letters) is a \textbf{generic AVR with integrated PSS and OEL}, defined via the CODEGEN \texttt{.txt} description. It provides a clean, modular three-block structure: AVR, PSS, and OEL.

\textbf{AVR (transient gain reduction + exciter):}

\[
\Delta V = V_0 - V + V_{PSS} - V_{OEL}
\]

\[
V_1(s) = G \cdot \frac{1 + s T_a}{1 + s T_b} \cdot \Delta V(s)
\]

\[
V_f(s) = \frac{1}{1 + s T_e} \cdot V_1(s), \quad V_f \in [V_{f,min},\, V_{f,max}]
\]

where \(V_0 = V + V_f/G\) is computed at initialisation.

\textbf{PSS (washout + two lead--lag stages):}

\[
V_{PSS}(s) = \frac{K_{PSS} \cdot s T_W}{1 + s T_W} \cdot \frac{1 + s T_1}{1 + s T_2} \cdot \frac{1 + s T_3}{1 + s T_4} \cdot \omega(s)
\]

Output is limited to \([-C,\, +C]\).

\textbf{OEL (inverse-time integrator):}

\[
x_{oel1} = I_f - if_{1lim}
\]

\[
\frac{d x_{oel2}}{dt} = \frac{1}{T_{oel}} \cdot \min(x_{oel1},\, if_{2lim}), \quad x_{oel2} \in [L_1,\, L_2]
\]

\[
V_{OEL} = \text{clip}(K_{oel} \cdot x_{oel2},\, 0,\, L_3)
\]

The CODEGEN block structure from the \texttt{.txt} file:

\begin{small}\begin{longtable}[]{@{}
  >{\raggedright\arraybackslash}p{(\linewidth - 6\tabcolsep) * \real{0.3529}}
  >{\raggedright\arraybackslash}p{(\linewidth - 6\tabcolsep) * \real{0.2321}}
  >{\raggedright\arraybackslash}p{(\linewidth - 6\tabcolsep) * \real{0.4149}}
@{}}
\toprule\noalign{}
\begin{minipage}[b]{\linewidth}\raggedright
Block
\end{minipage} & \begin{minipage}[b]{\linewidth}\raggedright
CODEGEN type
\end{minipage} & \begin{minipage}[b]{\linewidth}\raggedright
Signal
\end{minipage} \\
\midrule\noalign{}
\endhead
\bottomrule\noalign{}
\endlastfoot
AVR summing junction & \texttt{algeq} & \(V_0 - V + dv_{PSS} - dv_{OEL} - \Delta V\) \\
Transient gain reduction & \texttt{tf1p1z} & \(G(1+sT_a)/(1+sT_b)\) \\
Exciter & \texttt{tf1plim} & \(1/(1+sT_e)\) with \([V_{f,min}, V_{f,max}]\) \\
PSS washout & \texttt{tfder1p} & \(K_{PSS} \cdot sT_W / (1+sT_W)\) \\
PSS lead--lag 1 & \texttt{tf1p1z} & \((1+sT_1)/(1+sT_2)\) \\
PSS lead--lag 2 & \texttt{tf1p1z} & \((1+sT_3)/(1+sT_4)\) \\
PSS limiter & \texttt{lim} & \([-C, +C]\) \\
OEL error & \texttt{algeq} & \(I_f - if_{1lim}\) \\
OEL upper bound & \texttt{min1v1c} & \(\min(\cdot,\, if_{2lim})\) \\
OEL integrator & \texttt{inlim} & \(1/T_{oel}\) with \([L_1, L_2]\) \\
OEL multiplier & \texttt{algeq} & \(K_{oel} \cdot x_{oel2}\) \\
OEL output limiter & \texttt{lim} & \([0, L_3]\) \\
\end{longtable}\end{small}

\subsection{Parameters}\label{model-exc-custom:parameters-3}

\begin{small}\begin{longtable}[]{@{}ll@{}}
\toprule\noalign{}
Parameter & Description \\
\midrule\noalign{}
\endhead
\bottomrule\noalign{}
\endlastfoot
\texttt{G} & AVR gain (pu/pu) \\
\texttt{Ta} & AVR lead time constant / TGR zero (s) \\
\texttt{Tb} & AVR lag time constant / TGR pole (s) \\
\texttt{Te} & Exciter time constant (s) \\
\texttt{vfmin} & Minimum field voltage (pu) \\
\texttt{vfmax} & Maximum field voltage (pu) \\
\texttt{Kpss} & PSS gain \\
\texttt{Tw} & PSS washout time constant (s) \\
\texttt{T1} & PSS first lead--lag zero (s) \\
\texttt{T2} & PSS first lead--lag pole (s) \\
\texttt{T3} & PSS second lead--lag zero (s) \\
\texttt{T4} & PSS second lead--lag pole (s) \\
\texttt{C} & PSS output symmetric limit (pu) \\
\texttt{if1lim} & OEL field current threshold (pu) \\
\texttt{if2lim} & OEL maximum overcurrent error (pu) \\
\texttt{Toel} & OEL integrator time constant (s) \\
\texttt{Koel} & OEL output gain \\
\texttt{L1} & OEL integrator lower limit \\
\texttt{L2} & OEL integrator upper limit \\
\texttt{L3} & OEL output upper limit (pu) \\
\end{longtable}\end{small}

\textbf{Computed at initialisation:} \texttt{Vo\ =\ V\ +\ Vf\ /\ G}

\subsection{Usage Example}\label{model-exc-custom:usage-example-6}

\textbf{RAMSES Data File}

\begin{Verbatim}
SYNC_MACH GEN1  BUS1  1. 1.  0.  0.   600. 570. 3. 0. 0.95
                    XT   0.15  1.1  0.25  0.2  0.7  *  0.2  0.1  6.0257  0.  5.00  0.05  *  0.1
              EXC exc_GENERIC  200.0  0.05  10.0  0.04  0.0  5.0
                            ! G     Ta    Tb    Te    vfmin vfmax
                            9.5  1.40  0.154  0.033  0.154  0.033  0.10
                            ! Kpss Tw   T1     T2     T3     T4     C
                            1.05  0.30  30.0  1.0  0.0  1.0  2.0
                            ! if1lim if2lim Toel Koel L1  L2  L3
              TOR CONSTANT ;
\end{Verbatim}

\section{See Also}\label{model-exc-custom:see-also}

\begin{itemize}
\tightlist
\item
  \hyperref[doc:model-index]{Model Reference}, the full model catalogue
\item
  \hyperref[doc:model-exc-ieee]{IEEE Exciters}, the IEEE-standard excitation models
\item
  \hyperref[ch-sync]{Dynamic Data Records}, the \texttt{SYNC\_MACH} record these models sit in
\end{itemize}
